\subsection{Physical contractions and decoder allocation}

The physical screens operate on explicit literal tables. They try common
replacement actions for two controls, or separately redirected incoming
edges of an eliminated control. The replacement must reproduce the
specified tape, target control and head position inside the tested window,
and must avoid circular reliance on rows being rewritten. Passing one
representative window would not establish a global proof; all required
abstract contexts must pass, followed by the complete source and startup
checks. No proposed contraction passed the screens below.

\begin{center}\small
\begin{tabularx}{\linewidth}{@{}p{.24\linewidth}p{.14\linewidth}p{.18\linewidth}X@{}}
\toprule Input / rule & Radius & Step bound & Context scope\\\midrule
Primary PA194, pair merger & 4 & 80 & All-tape abstract contexts.\\
Primary PA194, pair merger & 6 & 120 & Blank-start abstract contexts.\\
Top-only Shift, pair merger & 6 & 300 & Blank-start abstract contexts; this input has 195 controls.\\
PA194DHS, pair merger & 7 & 300 & Blank-start abstract contexts; all working endpoints considered.\\
PA194DHS, incoming-edge bypass & 7 & 300 & All-tape abstract contexts; exact old two-step endpoint.\\
PA194DHS, joining-prefix pair & 4 & 40 / 120 & Blank-start abstract contexts; a later old configuration may be the common endpoint.\\
PA194DHS, joining-prefix pair & 6 & 80 / 240 & 274,589 abstract state/window pairs.\\
Top-only Shift, joining-prefix pair & 6 & 80 / 240 & 268,907 abstract state/window pairs.\\\bottomrule
\end{tabularx}\end{center}
Here a radius-$r$ window has $2r+1$ cells. These abstractions and bounded
continuations give sufficient local rewrite rules, not complete searches
through halting-equivalent machines. A failed overapproximation can hide
a reduction that requires a stronger global invariant.
For a joining-prefix screen, the two bounds concern the old prefix and
replacement continuation, respectively. Both paths must avoid rewritten
controls before their common endpoint. These three additional screens
also returned zero candidates.

A separate 15,120-case finite screen permutes all seven generated head
blocks in the DHS association and independently chooses the diagonal
argument from $A,D,H$. Its best allocation score is 194. Source semantics,
structured returns and simultaneous local folds still require their own
checks whenever an allocation score looks promising.

Three seeded grammar searches, including folds with eligible fixed
passive readers, likewise supplied no certified sub-194 table:
\begin{center}\small
\begin{tabularx}{\linewidth}{@{}Xp{.12\linewidth}r r@{}}
\toprule Search & Seconds & Proposal slots & Best score\\\midrule
DHS shared-subtree reassociation & 600 & 11,362,240 & 194\\
DHS head order, parameters and reassociation & 600 & 8,606,089 & 194\\
Normalized top-only Shift reassociation & 120 & 2,243,386 & 195\\\bottomrule
\end{tabularx}\end{center}
The seeds are 202611141931, 202611141932 and 202611141933, respectively.
The total of 22,211,715 counts proposal slots, including skipped mutations,
repeated trees and rejected proposals; it is not a count of distinct
machines or proofs. The corresponding accepted-mutation counts are
547,685, 3,012,683 and 130,995. The optional bundle at
\file{search/local_and_grammar/} records exact inputs, seeds, stopping
budgets, observations and limitations. Discovery is not required by
\file{verify_all.py}.

An early top-only mutation experiment had inconsistent indexing between
opaque and expanded increment macros. That run is excluded from trusted
search accounting. A corrected normalized reassociation experiment is
recorded separately. The accepted PA194DHS construction and its complete
verification do not depend on the rejected search run.
